\chapter*{Resumen}
\addcontentsline{toc}{chapter}{Resumen}
\markboth{Resumen}{Resumen}

El software comercial de análisis por el Método de los Elementos Finitos (MEF) opera como una caja negra: entrega resultados sin mostrar el procedimiento que los produce, de modo que el análisis no puede seguirse hasta sus datos ni comprobarse paso a paso. Esta tesis desarrolla EduFEM, un software educativo libre y en español para problemas bidimensionales de elasticidad lineal, en tensión y en deformación plana, con elementos cuadriláteros de cuatro y de nueve nodos (Q4 y Q9). La hipótesis sostiene que su desarrollo permitirá mejorar la forma en que el software expone el procedimiento de cálculo, elevando la trazabilidad y la verificabilidad del análisis. EduFEM reúne un motor de cálculo, una interfaz de pre-proceso, proceso y post-proceso sobre un mismo lienzo, ocho módulos educativos y una memoria de cálculo automática que puede reproducirse a mano. La trazabilidad se comprobó por cobertura: las nueve etapas del procedimiento quedan expuestas, y los dieciocho contenidos que un consenso publicado de expertos exige tienen instrumento en el software, según el cotejo del autor. La verificabilidad se comprobó con tres casos de estudio. Con soluciones manufacturadas, al reducir a la mitad el tamaño de los elementos, el error del desplazamiento se divide por cuatro con el Q4 y por ocho con el Q9, como predice la teoría. En la viga de Timoshenko, los errores frente a la solución analítica fueron del 0,04\,\% en la tensión normal y del 0,26\,\% en la flecha, y la diferencia frente a SAP2000, del 0,21\,\% en la tensión normal. La membrana de Cook mostró el bloqueo por cortante del Q4 frente a la convergencia del Q9. La hipótesis queda comprobada, dentro del alcance declarado, con criterios fijados de antemano.

\noindent\textbf{Palabras clave:} método de elementos finitos, software educativo, trazabilidad, verificación y validación, memoria de cálculo
